% !TEX root = ../b-rg.tex
\chapter{Conventions and Abbreviations}\label{ch:conventions}

\section*{Notation}
Borlish forms are given in italics, as in \bl{Borallesc}. Letters and spellings are given between angle brackets, as in \orth{ç}; phonemic transcriptions between slashes, as in \phm{sbɔk}; and phonetic transcriptions between square brackets, as in \pht{zbɔk}. Obsolete forms and senses are marked with a dagger\obs. Unless otherwise stated, phonetic transcriptions reflect the mesolectal speech of Damvath.

\section*{Examples}
Numbered examples give the Borlish form, its phonemic transcription (and, where relevant, its phonetic transcription), and a gloss. Glosses follow the Leipzig Glossing Rules: a hyphen separates affixes, an equals sign separates clitics, and a full stop joins several glosses that correspond to a single form. Grammatical categories are given in small capitals, using the abbreviations below. A category is glossed only where it is marked by a difference in form: the bare stem of a verb is glossed by its meaning alone, and a pronoun is glossed \Obl{} only where its oblique form differs from its subject form. For clarity, however, pronouns in disjunctive use are always glossed \Dsj{}, and third-person singular pronouns used as objects are always glossed \Acc{}.

\printnoidxglossary[type=\leipzigtype,style=long,nonumberlist,title={Glossing abbreviations}]
